\documentclass[12pt,a4paper]{article}

% ------------------------------------------------
% PACOTES
% ------------------------------------------------
\usepackage[utf8]{inputenc}
\usepackage[T1]{fontenc}
\usepackage[brazil]{babel}
\usepackage{geometry}
\usepackage{setspace}
\usepackage{hyperref}
\usepackage{graphicx}

% ------------------------------------------------
% CONFIGURACOES DE MARGEM E LAYOUT
% ------------------------------------------------
\geometry{
  left=3cm,
  right=2cm,
  top=3cm,
  bottom=2.5cm
}

\setstretch{1.5}

% ------------------------------------------------
\begin{document}

\begin{center}
    \Large\textbf{MANUAL DO PROFESSOR}\\[0.3cm]
    \large\textbf{Sistema de Registro de Presença por QR Code Dinâmico}\\[0.3cm]
    \normalsize Projeto-Piloto -- Faculdade de Engenharia Civil -- UNICAMP\\
    Disciplina: \textbf{Fundações}
\end{center}

\vspace{1cm}

% ------------------------------------------------

\section*{1. Apresentação}

Este manual descreve os procedimentos para utilização do sistema de registro automatizado de presença em sala de aula por meio de QR Code dinâmico, implantado como projeto-piloto institucional na disciplina \textbf{Fundações} da Faculdade de Engenharia Civil da UNICAMP.

O sistema foi concebido para otimizar o controle acadêmico de frequência, aumentando a agilidade do processo de chamada, reduzindo erros de registro e assegurando confiabilidade dos dados, em conformidade com a Lei Geral de Proteção de Dados (LGPD).

% ------------------------------------------------

\section*{2. Objetivos do Sistema}

O projeto visa:

\begin{itemize}
    \item Reduzir o tempo dedicado à chamada presencial;
    \item Minimizar erros manuais de registro;
    \item Aumentar a integridade dos dados de frequência;
    \item Facilitar auditorias acadêmicas;
    \item Gerar bases digitais para análise educacional.
\end{itemize}

% ------------------------------------------------

\section*{3. Funcionamento Geral}

O sistema utiliza QR Code temporário, gerados exclusivamente para cada aula. O fluxo operacional é descrito abaixo:

\begin{enumerate}
    \item O professor gera o QR Code ao iniciar a aula;
    \item O código é exibido em projetor ou tela;
    \item O aluno escaneia o QR com o celular;
    \item Informa RA e nome completo;
    \item O sistema registra automaticamente a presença em planilha vinculada à disciplina.
\end{enumerate}

Cada QR Code possui validade máxima de 3 minutos e torna-se automaticamente inválido após o vencimento.

% ------------------------------------------------

\section*{4. Segurança do Sistema}

As seguintes camadas de segurança estão ativas:

\begin{itemize}
    \item Geração de token único por aula;
    \item Expiração temporizada (3 minutos);
    \item Bloqueio de registros duplicados por RA;
    \item Registro automático de data e hora.
\end{itemize}

Estas medidas impedem o uso indevido do sistema fora do contexto presencial da aula.

% ------------------------------------------------

\section*{5. Conformidade com a LGPD}

São coletados exclusivamente:

\begin{itemize}
    \item Nome completo do aluno;
    \item Número de RA;
    \item Data e horário do registro.
\end{itemize}

O sistema não coleta dados biométricos, de localização ou imagens. A finalidade é restritamente acadêmica.

% ------------------------------------------------

\section*{6. Procedimento em Sala}

\subsection*{6.1. Geração do QR Code}

O professor acessa o endereço público do sistema e projeta o QR Code exibido.

\subsection*{6.2. Orientação aos Alunos}

Instruir os alunos a:

\begin{itemize}
    \item Escanear o QR Code com o celular;
    \item Informar RA e nome completo;
    \item Confirmar o envio.
\end{itemize}

\subsection*{6.3. Encerramento}

Após o período de validade do QR Code, novos registros são automaticamente bloqueados.

% ------------------------------------------------

\section*{7. Acompanhamento}

As presenças ficam disponíveis em planilha Google vinculada à disciplina:

\begin{center}
    Presenca\_FEC\_Fundacoes\_2025
\end{center}

Campos registrados:

\begin{itemize}
    \item Data
    \item Disciplina
    \item Turma
    \item RA
    \item Nome
    \item Horário
\end{itemize}

% ------------------------------------------------

\section*{8. Situações Especiais}

\begin{itemize}
    \item \textbf{Aluno sem celular:} registro manual posterior pelo professor.
    \item \textbf{Aluno atrasado:} sujeito à decisão do docente quanto ao lançamento manual.
    \item \textbf{Falhas técnicas:} utilizar lista manual e registrar posteriormente.
\end{itemize}

% ------------------------------------------------

\section*{9. Suporte Técnico}

Coordenação técnica do projeto:

\begin{center}
    \textbf{Heloi Moacyr Tanoue}
\end{center}

% ------------------------------------------------

\section*{10. Avaliação do Piloto}

Serão analisados:

\begin{itemize}
    \item Taxa de adesão dos alunos;
    \item Redução no tempo de chamada;
    \item Incidência de erros ou bloqueios;
    \item Feedback docente.
\end{itemize}

% ------------------------------------------------

\section*{11. Expansão}

Em caso de validação positiva do piloto, recomenda-se a ampliação para outras disciplinas da FEC e integração futura aos sistemas institucionais.

% ------------------------------------------------

\vfill
\begin{center}
    Versão 1.0 -- Novembro de 2025
\end{center}

\end{document}

